\section{Variaciones concretas de Flow Matching}

\subsection{Opción más simple: trayectorias de Transporte Óptimo}

La opción más sencilla y frecuentemente utilizada es el \textbf{mapeo de flujo condicional de Transporte Óptimo (OT)}, también conocido como rectified flow:

$$\phi_t(x_0 | x_1) = (1 - t) x_0 + t \, x_1$$

Los parámetros correspondientes son:
\begin{itemize}
    \item $\mu_t(x_1) = t \, x_1$
    \item $\sigma_t(x_1) = 1 - t$
\end{itemize}

Esto constituye una \textbf{interpolación lineal} entre $x_0$ y $x_1$, donde la trayectoria condicional correspondiente es una \textbf{línea recta}.

\begin{importantbox}{Campo vectorial y objetivo de entrenamiento de la ruta OT}
Para la ruta OT, el campo vectorial condicional es extremadamente simple:
$$u_t(x | x_1) = x_1 - x_0$$
Este es un \textbf{vector constante} (no depende de $t$), con dirección desde $x_0$ hacia $x_1$.

Por lo tanto, el objetivo de entrenamiento de Conditional Flow Matching se convierte en:
$$\mathcal{L}_{\text{CFM}}(\theta) = \mathbb{E}_{t, \, x_1, \, x_0} \left\| v_\theta(t, (1-t)x_0 + t x_1) - (x_1 - x_0) \right\|^2$$

El proceso de entrenamiento es extremadamente sencillo:
\begin{enumerate}
    \item Muestrear $x_1 \sim p_{\text{data}}$, $x_0 \sim \mathcal{N}(0, I)$, $t \sim \mathcal{U}[0,1]$
    \item Calcular $x_t = (1-t)x_0 + t \, x_1$
    \item Entrenar la red $v_\theta(t, x_t)$ para predecir $x_1 - x_0$
\end{enumerate}
\end{importantbox}

\subsection{Correspondencia con los schedulers VP/VE en modelos de difusión}

Además de la ruta OT, también podemos elegir una parametrización que corresponda a los schedulers VP (Variance Preserving) o VE (Variance Exploding) de los modelos de difusión.

\textbf{Correspondencia VP} (análogo al scheduler de ruido de DDPM):
\begin{itemize}
    \item $\mu_t(x_1) = \alpha_t \, x_1$, donde $\alpha_0 = 0$, $\alpha_1 \approx 1$
    \item $\sigma_t(x_1) = \sqrt{1 - \alpha_t^2}$ (manteniendo la varianza total constante)
\end{itemize}

El campo vectorial condicional se convierte en:
$$u_t(x | x_1) = \frac{\alpha_t'}{1 - \alpha_t^2} \left( \alpha_t \, x_1 - x \right) + \alpha_t' \, x_1$$

\begin{knowledgebox}{Relación entre Flow Matching y la función de puntuación (score function)}
Existe una relación matemática exacta entre el campo vectorial condicional y la función de puntuación (\textit{score function}) en los modelos de difusión. Bajo la parametrización VP:
$$u_t(x | x_1) = f(t) \, x + g(t) \, \nabla_x \log p_t(x | x_1)$$
donde $f(t)$ y $g(t)$ son coeficientes relacionados con el scheduler de ruido.

Esto indica que \textbf{entrenar el campo vectorial de Flow Matching y la función de puntuación para modelos de difusión están esencialmente haciendo lo mismo}—solo utilizan diferentes parametrizaciones y pesos de pérdida distintos.
\end{knowledgebox}

\subsection{Trayectorias rectas vs trayectorias curvas}

\begin{warningbox}{Trayectoria condicional vs trayectoria marginal}
Aunque cada \textbf{trayectoria condicional} (dada $x_0$ y $x_1$) es una línea recta, la \textbf{trayectoria marginal} (solo dada $x_0$, sin conocer el objetivo $x_1$) generalmente no lo es. Esto se debe a que un mismo punto espacial $x_t$ puede ser atravesado por múltiples trayectorias condicionales distintas, y la red debe predecir el promedio ponderado de estos campos vectoriales condicionales.

Esta es la razón por la cual, incluso al entrenar con la ruta OT, sigue siendo necesario resolver una ODE en múltiples pasos durante el muestreo real—ya que el campo vectorial marginal aprendido no es constante.
\end{warningbox}

\subsection{Resumen de la sección}

La ruta OT (interpolación lineal) es la variación más simple y frecuentemente utilizada de Flow Matching, donde el objetivo de entrenamiento se reduce a predecir $x_1 - x_0$. Al seleccionar diferentes funciones para $\mu_t$ y $\sigma_t$, Flow Matching puede reproducir diversos schedulers de ruido de los modelos de difusión. Esta flexibilidad constituye una ventaja importante de los modelos de flujo.
